\documentclass[11pt,a4paper]{ctexart}
\usepackage[margin=22mm]{geometry}
\usepackage{amsmath,amssymb,mathtools,booktabs,longtable,array,enumitem,fancyhdr,xcolor}
\usepackage[colorlinks=true,linkcolor=blue!50!black,urlcolor=blue!50!black]{hyperref}
\setlength{\headheight}{15pt}
\pagestyle{fancy}\fancyhf{}\fancyhead[L]{微分方程 I · 期中复习}\fancyhead[R]{截止 3.1 节（含）}\fancyfoot[C]{\thepage}
\setlength{\parindent}{0pt}\setlength{\parskip}{5pt}
\setlist{itemsep=3pt,topsep=4pt,leftmargin=2em}
\newcommand{\R}{\mathbb R}
\newcommand{\dd}{\,\mathrm d}
\newcommand{\tr}{\operatorname{tr}}
\newcommand{\rank}{\operatorname{rank}}
\newcommand{\point}[2]{\subsubsection*{#1}\label{kp:#2}}
\newcommand{\core}{\textbf{【重点】}}
\allowdisplaybreaks[2]
\begin{document}

\title{微分方程 I：期中考试复习重点\\\large 考试截止到第 3.1 节（含）}
\author{}\date{2026 年 10 月 31 日修订}
\maketitle\vspace{-1.5em}
\textbf{考试范围。}《微分方程 I（101 计划）》第 1、2 章及第 3.1 节（含）；教材印刷页 1--94（PDF 页 11--104，第 94 页下半的 3.2 除外）。\textbf{不含第 3.2 节及以后内容。}

\textbf{材料依据。}教材、7 组作业材料（8 份 PDF，修订稿去重后计 7 组，其中 1 组属 3.2，排除），以及所提供的两份历年试题。作业映射与试卷题型见附录；历年分值不代表本次考试分配。

\textbf{约定。}默认实数范围；非初等积分可保留为积分形式。训练题及出处另册列出。
\setcounter{tocdepth}{1}\tableofcontents\newpage

\section{第一章：基本概念}
\setcounter{subsection}{1}
\subsection{微分方程概念}
\point{2. 阶、线性与标准形式}{definition}
$F(t,x,x',\ldots,x^{(n)})=0$ 为一般隐式形式，最高出现的导数阶数为方程的阶；导数的幂不改变阶数。若能解出最高阶导数，则为显式形式 $x^{(n)}=f(t,x,\ldots,x^{(n-1)})$。$F_{x^{(n)}}\ne0$ 是局部解出最高阶导数的常用充分条件。

线性方程可写为
\[
a_n(t)x^{(n)}+\cdots+a_1(t)x'+a_0(t)x=b(t).
\]
系数 $a_j,b$ 仅依赖 $t$，未知函数及其各阶导数均线性出现。$b\equiv0$ 时为齐次方程；$b\not\equiv0$ 时为非齐次方程。在 $a_n(t)\ne0$ 的区间内可化为规范形式；$a_n(t)=0$ 处单独讨论。

\point{3. 解、隐式解、通解、特解与奇解}{general}
设 $I$ 为区间。若 $x\in C^n(I)$ 的图像位于方程定义域内，且
\[
F(t,x(t),x'(t),\ldots,x^{(n)}(t))=0\quad(t\in I),
\]
则称 $x$ 为该方程在 $I$ 上的解。隐式关系 $U(t,x)=0$ 所确定的可微分支满足原方程时称为隐式解；$U_x\ne0$ 为局部解出 $x(t)$ 的充分条件。

$n$ 阶方程的通解 $x=\phi(t,c_1,\ldots,c_n)$ 含 $n$ 个独立常数，教材的独立性判据为
\[
\det\frac{\partial(\phi,\partial_t\phi,\ldots,\partial_t^{n-1}\phi)}{\partial(c_1,\ldots,c_n)}\ne0.
\]
上述函数族须满足原方程。通解的隐式表达称通积分；确定参数所得的解称特解。\textbf{通解不必包含全部解。}奇解是具有下述性质的解：在其每一点的任意邻域内，存在另一条与其在该点相切的解曲线。

\textbf{包络（教材定义）。}设 $K(C):U(t,x,C)=0$，$\gamma$ 为 $C^1$ 曲线。若对任意 $q\in\gamma$，存在 $C_*$，使 $K(C_*)$ 与 $\gamma$ 在 $q$ 处相切且在 $q$ 的某邻域内不重合，则称 $\gamma$ 为该曲线族的一支包络。
\[
U(t,x,C)=0,\qquad U_C(t,x,C)=0
\]
消去 $C$ 所得曲线为包络候选；须验证包络条件。若用其求奇解，还须验证原微分方程。

\point{4. 初值、边值、存在区间}{ivp}
$n$ 阶方程的初值条件为
\[
x^{(j)}(t_0)=a_j,\qquad j=0,\ldots,n-1.
\]
边值条件在不同端点给定。隐式方程的初值须与方程相容。初值解的存在区间须包含 $t_0$，并满足原方程定义域、实幂与对数的取值条件及所需可微性。

\point{5. 由曲线族消去常数建立微分方程}{eliminate}
对含 $n$ 个独立参数的函数族，联立原式及其前 $n$ 阶导数，消去参数，构造 $n$ 阶微分方程；隐式族可隐式求导后消元。除法、平方与开方所得关系须代回验证。


\textbf{例：由两参数通解构造二阶方程。}
题设称
\[
y=\sqrt{c_1x^2+c_2}
\]
为某微分方程的通解。参数独立时，所求方程应为二阶。
由 $y^2=c_1x^2+c_2$ 连续求导，得
\[
yy'=c_1x,\qquad (y')^2+yy''=c_1.
\]
消去 $c_1$ 得到
\[
\boxed{x\bigl((y')^2+yy''\bigr)-yy'=0.}
\]
在 $x\ne0$、$y>0$ 的局部区间上，令 $u=y^2$，上述方程等价于 $xu''-u'=0$，其通解为 $u=C_1x^2+C_2$，取正平方根即恢复题给函数族。此时
\[
\det\frac{\partial(y,y')}{\partial(c_1,c_2)}=-\frac{x}{2y^2}\ne0,
\]
故参数独立。$y=0$ 或 $x=0$ 处另行检验。该两参数族即使满足某三阶方程，也不构成该三阶方程的通解。

\subsection{微分方程形式}
\point{6. 高阶方程与混合阶系统的一阶化}{firstorder}
对 $x^{(n)}=f(t,x,\ldots,x^{(n-1)})$ 令 $u_i=x^{(i-1)}$，则
\[
u_1'=u_2,\ldots,u_{n-1}'=u_n,\qquad u_n'=f(t,u_1,\ldots,u_n).
\]
初值同步转成向量初值。混合阶系统对每个未知函数保留到其最高阶的前一阶导数，维数为各个已解出最高阶导数的阶数之和。例如 $x''+y=0,y''-x=0$ 化成 $(x,x',y,y')'=(x',-y,y',x)$。

\point{7. 对称形式、方向场与等倾线}{direction}
$M(t,x)\dd t+N(t,x)\dd x=0$ 在 $N\ne0$ 处可写成 $x'=-M/N$，在 $M\ne0$ 处可写成 $t'=-N/M$。对称形式可描述竖直切线，显式 $x(t)$ 则不能把竖直线当作函数解。

$x'=f(t,x)$ 的积分曲线在 $(t_0,x_0)$ 处的切线为 $x-x_0=f(t_0,x_0)(t-t_0)$。等倾线为 $f(t,x)=k$；其本身未必为积分曲线。自治标量方程的平衡解满足
\[
x'=g(x),\qquad x(t)\equiv a\iff g(a)=0.
\]

\section{第二章：初等积分法}
\subsection{变量分离的微分方程}
\point{9. 分离变量的标准结果}{separate}
若 $x'=a(t)b(x)$，$a,b$ 连续，则在 $b(x)\ne0$ 的连通区间上，令 $B'=1/b,A'=a$，有
\[
B(x)=A(t)+C,\qquad x=B^{-1}(A(t)+C).
\]
$B$ 在该区间严格单调；只有 $A(t)+C\in B(J)$ 时反函数才有意义。$b(a)=0$ 给出恒等解 $x\equiv a$。初值 $b(x_0)\ne0$ 时可写 $\int_{x_0}^{x(t)}\!\dd u/b(u)=\int_{t_0}^t a(s)\dd s$。

\point{10. 除法漏解、零解拼接与分支限制}{glue}
\core 在分离变量的过程中，除以 $b(x)$ 将讨论范围限制为 $b(x)\ne0$。若 $b(a)=0$，则常数函数 $x(t)\equiv a$ 满足原方程，但不属于该除法变形后的方程，因而应将这些常数解与分离变量所得的非平衡解共同列出。

平方或开方后得到的显式表达式仍须满足变形前的符号条件。例如 $x'=\sqrt{x}$ 的非零解支为 $x=(t-a)^2/4$、$t>a$；在 $t<a$ 时，该表达式的导数为负，不满足原方程。与零解连接后得到
\[
x(t)=\begin{cases}0,&t\le a,\\(t-a)^2/4,&t>a,\end{cases}
\]
该函数在整个实轴上为 $C^1$，并逐点满足原方程。

对作业中的 $x'=x^{1/3},x(0)=0$，除恒零解外，对任意 $a\ge0$ 都有
\[
x(t)=\begin{cases}0,&t\le a,\\ \displaystyle\pm[2(t-a)/3]^{3/2},&t>a.\end{cases}
\]
这说明连续右端不保证唯一性。对 $x'=g(x)$ 在非零一侧，到达零点 $a$ 的时间由 $\int\!\dd x/g(x)$ 是否在该端点收敛判断；收敛时才可能有限时刻接到平衡解，拼接后仍需验证可微性与原方程。

\point{11. Logistic 解、平衡解与渐近行为}{logistic}
对 $x'=kx(1-x)$、$k>0$、$x(t_0)=x_0$，有
\[
x(t)=\frac{x_0e^{k(t-t_0)}}{1-x_0+x_0e^{k(t-t_0)}}.
\]
该式也涵盖 $x_0=0,1$。$0<x_0<1$ 时在 $\R$ 上存在，从 $0$ 单调趋向 $1$；$x_0>1$ 时向前递减趋于 $1$，向后有有限极点；$x_0<0$ 时向前有有限极点。极点为 $t_*=t_0+k^{-1}\ln((x_0-1)/x_0)$，取不跨越它且含 $t_0$ 的区间。一般承载量 $K$ 的式子由 $x/K$ 归一化得到。

\point{12. 第一类：零次齐次方程}{homogeneous}
\core 以下第 12--14 条分类列出可化为变量分离方程的主要类型。式中 $x=x(t)$，$F,g$ 为给定函数；换元只在原式及新变量均有定义的区间上使用。

\textbf{类型：}$x'=g(x/t)$，或右端 $f(t,x)$ 满足 $f(\lambda t,\lambda x)=f(t,x)$。令 $v=x/t$（$t\ne0$），则
\[
x'=v+tv',\quad tv'=g(v)-v,\quad
\int\frac{\dd v}{g(v)-v}=\ln|t|+C.
\]
若 $g(v_0)=v_0$，则 $v(t)\equiv v_0$ 满足变换后的方程，故 $x(t)=v_0t$ 是原方程的解。在分离变量时除以 $g(v)-v$ 会排除这些常值 $v$，因此完整的求解结果应同时列出上述直线解。此处的零次齐次指右端的缩放性质，与线性齐次方程的定义不同。

\textbf{例：}$x'=x/t+(x/t)^2$。换元后 $tv'=v^2$，非零支满足 $-1/v=\ln|t|+C$，即
\[
x(t)=-\frac{t}{\ln|t|+C}.
\]
此外，$v\equiv0$ 对应原方程的解 $x\equiv0$，它不能由上式取有限常数 $C$ 得到。非零解与零解均应列入求解结果；各解的存在区间不得跨越 $t=0$，非零解还须满足 $\ln|t|+C\ne0$。

\point{13. 第二、三类：两个一次式之比}{affine-ratio}
统一类型为
\[
x'=F\!\left(\frac{at+bx+c}{mt+nx+d}\right),\qquad mt+nx+d\ne0.
\]
\textbf{第二类：$\Delta=an-bm\ne0$，先平移再作比值换元。}解
\[
at_*+bx_*+c=0,\qquad mt_*+nx_*+d=0,
\]
令 $T=t-t_*,X=x-x_*$，则 $\dd X/\dd T=F((aT+bX)/(mT+nX))$。在 $T\ne0$ 上再令 $v=X/T$，得
\[
T\frac{\dd v}{\dd T}=F\!\left(\frac{a+bv}{m+nv}\right)-v.
\]
若 $v_0$ 满足 $m+nv_0\ne0$ 且 $F((a+bv_0)/(m+nv_0))=v_0$，则 $X=v_0T$ 满足变换后的方程，对应原坐标中的直线 $x=x_*+v_0(t-t_*)$。分离变量时被除去的这些直线解应一并列出，其有效区间仍由原分母非零条件确定。

\textbf{例（作业习题 2.1 第 1(3) 题）：}
\[
x'=\frac{2t+3x-1}{t-x+1}.
\]
交点为 $(-2/5,3/5)$。取 $T=t+2/5,X=x-3/5,v=X/T$，得到
\[
T v'=\frac{v^2+2v+2}{1-v},\qquad
\frac{1-v}{v^2+2v+2}\dd v=\frac{\dd T}{T}.
\]
这里 $v'$ 指对 $T$ 求导。原分母为零的点不属于原方程定义域；$T=0$ 只是本次比值换元的限制，原解能否跨过它须回原式判断。

\textbf{第三类：$\Delta=0$，用共同线性组合。}若 $b,n$ 不同时为零，则存在 $\lambda$ 使 $a=\lambda b,m=\lambda n$。令 $z=\lambda t+x$，得
\[
z'=\lambda+F\!\left(\frac{bz+c}{nz+d}\right),
\]
为自治、可分离方程。若 $b=n=0$，右端仅含 $t$，直接积分；若两个一次式恒成比例且比值有定义，右端为常数。这些退化情形不能强行使用交点平移法。

\textbf{例：}$x'=(t+x+1)/(2t+2x-1)$。取 $z=t+x$，得
\[
z'=\frac{3z}{2z-1},\qquad
\frac{2z-1}{3z}\dd z=\dd t.
\]
非零支满足 $2z-\ln|z|=3t+C$，即 $2(t+x)-\ln|t+x|=3t+C$；另有 $z\equiv0$ 对应的解 $x=-t$。始终要求 $2z-1\ne0$。

\point{14. 第四、五类：线性组合与加权齐次方程}{weighted}
\textbf{第四类：}$x'=F(at+bx+c)$。令 $z=at+bx+c$，则
\[
z'=a+bF(z).
\]
当 $b\ne0$ 时由 $x=(z-at-c)/b$ 回代；当 $b=0$ 时原方程直接对 $t$ 积分。$a+bF(z_0)=0$ 所对应的常值 $z$ 解不可遗漏。

\textbf{例（教材例 2.1.5）：}$x'=\sin(3t+2x)$。取 $z=3t+2x$，得 $z'=3+2\sin z$，故
\[
\int\frac{\dd z}{3+2\sin z}=t+C,\qquad z=3t+2x.
\]
由于 $3+2\sin z>0$，不存在常值 $z$ 解，此步除法不会漏解。作业中的 $x'=e^{-x+2t}+2$ 则取 $z=x-2t$，化成 $z'=e^{-z}$，得 $e^{x-2t}=t+C$。

\textbf{第五类（补充）：加权齐次方程。}在 $t>0$ 上，若
\[
f(\lambda^{s_1}t,\lambda^{s_2}x)=\lambda^k f(t,x),\quad
s_1\ne0,\quad k=s_2-s_1,
\]
则 $f(t,x)=t^{q-1}H(x/t^q)$，其中 $q=s_2/s_1,H(v)=f(1,v)$。令 $x=t^qv$，得
\[
tv'=H(v)-qv,\qquad\frac{\dd v}{H(v)-qv}=\frac{\dd t}{t}.
\]
若 $H(v_0)=qv_0$，则变换后的方程有常值解 $v\equiv v_0$，对应原方程的幂函数解 $x=v_0t^q$。这些解不由除以 $H(v)-qv$ 后的积分式直接给出，应与非平衡解共同列出。上述实幂公式以 $t>0$ 为前提；在 $t<0$ 上采用同类换元时，须重新指定实幂的定义及相应符号。

\textbf{例：}$x'=x/t+x^2/t^3$（$t>0$）。取 $x=t^2v$，得 $tv'=v(v-1)$，故非平衡支满足
\[
\ln\left|\frac{v-1}{v}\right|=\ln t+C,
\qquad x=\frac{t^2}{1-C_1t}.
\]
分离时除以 $v(v-1)$ 排除了 $v\equiv0$ 与 $v\equiv1$。前者对应解 $x\equiv0$，不能由上述显式公式取有限 $C_1$ 得到；后者对应解 $x=t^2$，已经由上式取 $C_1=0$ 给出。因此求解结果为上述函数族与零解，且始终要求 $t>0$、$1-C_1t\ne0$。

各类换元均须保留原定义域，并补充因除法遗漏的常值解。

\point{15. 隐式积分、曲线极限与渐近线}{asymptote}
设 $B:J\to B(J)$ 连续且严格单调，$B(x(t))=A(t)+C$。若
\[
\lim_{t\to\pm\infty}A(t)=L\in\R,\qquad L+C\in B(J),
\]
则 $\lim_{t\to\pm\infty}x(t)=B^{-1}(L+C)$，相应水平渐近线为 $x=B^{-1}(L+C)$。极限落在值域边界时，须另行分析反常积分与反函数极限。

\subsection{一阶线性微分方程}
\point{16. 积分因子与初值公式}{linear-scalar}
\core 设 $p,q$ 在区间 $I$ 连续。对标准形式
\[
\boxed{x'+p(t)x=q(t),}
\]
取 $p$ 的一个固定原函数 $P$，即 $P'=p$，则积分因子为 $\mu(t)=e^{P(t)}$，满足 $(\mu x)'=\mu q$。\textbf{通解公式}为
\[
\boxed{x(t)=Ce^{-P(t)}+e^{-P(t)}Q(t),\qquad Q'(t)=q(t)e^{P(t)}.}
\]
其中各处 $\int p$ 均取同一个原函数 $P$，$Q$ 为 $qe^P$ 的一个固定原函数；$C$ 为唯一任意常数。

\textbf{定积分形式（明确内外层变量与上下限）。}固定 $t_0\in I$，令 $P(t)=\int_{t_0}^t p(u)\dd u$，则
\[
\boxed{
x(t)=\exp\!\left(-\int_{t_0}^t p(u)\dd u\right)
\left[C+\int_{t_0}^t q(s)\exp\!\left(\int_{t_0}^s p(u)\dd u\right)\dd s\right].}
\]
\textbf{初值解公式：}若 $x(t_0)=x_0$，则 $C=x_0$，等价地
\[
\boxed{x(t)=\exp\!\left(-\int_{t_0}^t p(u)\dd u\right)
\left[x_0+\int_{t_0}^t q(s)\exp\!\left(\int_{t_0}^s p(u)\dd u\right)\dd s\right].}
\]
\textbf{任意积分起点。}对固定 $a,b\in I$，令 $P_a(t)=\int_a^t p(u)\dd u$，通解亦可写为
\[
x(t)=e^{-P_a(t)}\left[C+\int_b^t q(s)e^{P_a(s)}\dd s\right].
\]
初值 $x(t_0)=x_0$ 对应
\[
C=e^{P_a(t_0)}x_0-\int_b^{t_0}q(s)e^{P_a(s)}\dd s.
\]
取 $a=b=t_0$ 时，$P_a(t_0)=0$、$C=x_0$；改变起点时应按上式重新确定 $C$。初值解在整个 $I$ 上存在且唯一。齐次解为 $Ce^{-P}$，非零齐次解处处非零且符号不变。

\point{17. 解的叠加与两个已知解}{affine}
非齐次方程两解之差是齐次解，通解为“一个特解＋齐次通解”。若 $x_1,x_2$ 是不同的非齐次解，则
\[
x=x_1+C(x_2-x_1)=(1-C)x_1+Cx_2
\]
就是通解；二者在任何时刻都不同。任意两个非齐次解的和一般不是同一方程的解；常系数仿射组合的系数和为 $1$ 时仍为解。

\point{18. 复合未知量与交换自变量}{linearize}
若方程具有形式 $h'(x)x'+p(t)h(x)=q(t)$，令 $z=h(x)$，则由链式法则得到线性方程 $z'+p(t)z=q(t)$。例如取 $h(x)=\sin x$ 时，$z'=\cos x\,x'$。求得 $z(t)$ 后，应在 $h$ 的值域内求解 $h(x(t))=z(t)$，选择满足原方程与初值的可微实数分支。若推导中除以 $h'(x)$，则该推导只适用于 $h'(x)\ne0$ 的区域；$h'(x)=0$ 处的常值解及可能延拓须直接代入原方程检验。

若交换自变量后得到 $\dd t/\dd x+a(x)t=b(x)$，可先求线性方程的解 $t=t(x)$。只有在 $\dd t/\dd x\ne0$ 的局部区间内，才能由反函数定理恢复可微函数 $x=x(t)$；不满足该条件的点须根据原方程单独讨论。

\point{19. Bernoulli 标准变换与实幂分支}{bernoulli}
考虑方程
\[
x'+p(t)x=q(t)x^\alpha,\qquad p,q\in C(I).
\]
\textbf{标准换元。}设 $\alpha\ne0,1$，记 $\beta=1-\alpha\ne0$。在实幂有定义且 $x\ne0$ 的解支上，除以 $x^\alpha$，令 $z=x^\beta$，由链式法则得
\[
z'=\beta x^{-\alpha}x',\qquad
\boxed{z'+\beta p(t)z=\beta q(t).}
\]
对非零初值 $x(t_0)=x_0$，应取 $z(t_0)=x_0^\beta$。令 $P(t)=\int_{t_0}^t p(s)\dd s$，则
\[
z(t)=e^{-\beta P(t)}\left[x_0^\beta+
\beta\int_{t_0}^t q(s)e^{\beta P(s)}\dd s\right].
\]
在选定的正、负半轴上，$\dd z/\dd x=\beta x^{-\alpha}\ne0$，故该变换可逆。以下实幂定义域、逆分支及例子为\textbf{补充说明}。

\textbf{实幂定义域。}对任意实数 $\alpha$，$x>0$ 时均可定义 $x^\alpha=e^{\alpha\ln x}$。在 $x<0$ 上采用通常的实数幂约定时，只有 $\alpha$ 为整数，或为最简分数 $m/n$ 且 $n$ 为奇数，才能使用相应实幂；整数包括在 $n=1$ 的情形中。此时
\[
x^{m/n}=(-1)^m|x|^{m/n}\quad(x<0).
\]
若最简分母为偶数或 $\alpha$ 为无理数，则不能自行增加负解支。$\alpha>0$ 时 $0^\alpha=0$，原方程允许零解；$\alpha<0$ 时 $x=0$ 不属于原方程定义域，零函数不是解。

\textbf{逆变换的值域限制。}正解支始终要求
\[
\boxed{z>0,\qquad x=z^{1/\beta}.}
\]
若允许负解，写 $\alpha=m/n$ 为最简分数，$n$ 为奇数，则 $z=x^{(n-m)/n}=(-1)^{n-m}|x|^\beta$。因此
\[
\begin{array}{c|c|c}
\text{选定的原解支}&\text{允许的 }z&\text{逆变换}\\\hline
x>0&z>0&x=z^{1/\beta}\\
x<0,\ m\text{ 为奇数}&z>0&x=-z^{1/\beta}\\
x<0,\ m\text{ 为偶数}&z<0&x=-(-z)^{1/\beta}
\end{array}
\]
表中每个带实指数的底数均为正数。在给定初值时，应选取允许值域内包含 $t_0$ 的连通时间区间，并由 $x_0$ 的符号确定逆分支。线性方程的任意实解 $z(t)$ 不一定对应原方程的合法实解。

\textbf{典型指数。}下表均先讨论 $x\ne0$ 的解支；$\alpha>0$ 时零解须另外计入，$\alpha<0$ 时不得加入零解。
\[
\begin{array}{c|c|c|c}
\alpha&z=x^{1-\alpha}&\text{允许的 }z&\text{可选逆分支}\\\hline
\frac12&\sqrt x&z>0&x=z^2\\
\frac13&x^{2/3}=|x|^{2/3}&z>0&x=\pm z^{3/2}\\
\frac23&x^{1/3}&z\ne0&x=z^3\\
\frac32&x^{-1/2}&z>0&x=z^{-2}\\
2&x^{-1}&z\ne0&x=1/z\\
3&x^{-2}&z>0&x=\pm z^{-1/2}\\
-1&x^2&z>0&x=\pm\sqrt z
\end{array}
\]
例如 $\alpha=1/2$ 时，虽然负 $z$ 也可平方，但 $\sqrt{z^2}=|z|\ne z$，故不能用负 $z$ 恢复原换元。对 $x'=\sqrt x$，线性方程给出 $z=(t-a)/2$，所以 $x=(t-a)^2/4$ 仅在 $t>a$ 时为非零解支；$t<a$ 时该表达式的导数为负，不满足原方程。

类似地，$\alpha=3/2$ 时，$x=z^{-2}$ 虽在负 $z$ 上也有实数值，但 $x^{-1/2}=|z|\ne z$，仍须保留 $z>0$。$\alpha=1/3$ 与 $\alpha=3$ 则各有正、负两个逆分支，不能只保留正支。

\textbf{退化指数。}$\alpha=0$ 时按常数项 $x^0=1$ 处理，原方程为 $x'+px=q$；零函数为解当且仅当 $q\equiv0$。$\alpha=1$ 时原方程为 $x'+(p-q)x=0$，而 $z=x^0=1$ 不能区分不同的解，故不得使用上述可逆变换。

\point{20. Bernoulli 解能否经过零点}{bernoulli-zero}
\core 设 $p,q$ 连续，实幂取合法分支。$\alpha>1$ 时，右端在零附近局部 Lipschitz，故 $x(t_0)=0\Rightarrow x\equiv0$。$0<\alpha<1$ 时，零初值解可能不唯一；$q\equiv0$ 时仍唯一。

\textbf{$0<\alpha<1$ 时的零点连接。}此时 $0<\beta=1-\alpha<1$。在非负分支上，给定 $a\in I$，取线性方程满足 $z(a)=0$ 的解
\[
z(t)=\beta e^{-\beta\int_a^t p(u)\dd u}
\int_a^t q(s)e^{\beta\int_a^s p(u)\dd u}\dd s.
\]
若 $z$ 在 $a$ 的某一侧为正，则该侧取 $x=z^{1/\beta}$，另一侧取 $x=0$。由于
\[
x'=\frac1\beta z^{1/\beta-1}z'\longrightarrow0\quad(t\to a),
\]
该连接为原方程的局部 $C^1$ 解。若 $q(a)\ne0$，由 $z'(a)=\beta q(a)$ 可知正支存在于 $a$ 的某一侧；若 $q(a)=0$，则须另行检查 $z$ 的符号。在零解持续的时间段内，$z\equiv0$ 通常不满足该非齐次线性方程，所以须在各非零区间分别换元，并直接验证零点处的连接。

不同指数的变号性质并不相同。例如
\[
x'=x^{2/3}:\qquad x(t)=\frac{(t-a)^3}{27}
\]
是可从负值经过零点变为正值的全实轴解；而 $x'=x^{1/3}$ 的非零分支为 $\pm[2(t-a)/3]^{3/2}$、$t>a$，可以与此前的零解连接。

\textbf{其他指数的端点。}$\alpha>1$ 时 $\beta<0$，正分支上 $z\to0^+$ 对应 $x\to+\infty$，并非 $x\to0$；此时到达 $z=0$ 通常表示原解发散。$\alpha<0$ 时，$z\to0^+$ 虽对应 $x\to0^+$，但原方程在零点无定义，不能接上零解或作为原方程的解穿过零点。

\textbf{例：}$x'-4x/t=t^2\sqrt{x}$ 的非零解支满足
\[
z=\sqrt{x}=t^2(t/2+C),\qquad x=t^4(t/2+C)^2,
\qquad t\ne0,\quad t/2+C>0.
\]

\point{21. Riccati 已知一个、两个、三个解时的结果}{riccati}
设 $p,q,r$ 在区间 $I$ 上连续，考虑 Riccati 方程
\[
x'+p(t)x=q(t)x^2+r(t).
\]
\textbf{已知一个特解。}若已知解 $x_1$，在 $x\ne x_1$ 的区间上令 $u=1/(x-x_1)$，即 $x=x_1+1/u$，则
\[
u'+(2qx_1-p)u=-q.
\]
求得该线性方程的解 $u(t)$ 后，在 $u(t)\ne0$ 的区间上由 $x(t)=x_1(t)+1/u(t)$ 得到原方程的解。由于 $1/u$ 不可能取零，已知解 $x=x_1$ 不属于此换元得到的函数族；完整的求解结果还应明确列出 $x(t)=x_1(t)$。

\textbf{已知两个不同解。}若 $x_1,x_2$ 为 $I$ 上的两个不同解，固定 $t_*\in I$，定义
\[
E(t)=\exp\!\left(\int_{t_*}^t q(s)\bigl(x_1(s)-x_2(s)\bigr)\dd s\right).
\]
对满足 $x\ne x_2$ 的解，有
\[
\frac{x(t)-x_1(t)}{x(t)-x_2(t)}=CE(t),\qquad
\boxed{x(t)=\frac{x_1(t)-CE(t)x_2(t)}{1-CE(t)}.}
\]
其中 $C$ 为任意实常数，解的区间须满足 $1-CE(t)\ne0$。取 $C=0$ 时，上式给出已知解 $x=x_1$。已知解 $x=x_2$ 在构造比值时被分母排除，且不能由该显式公式取任何有限 $C$ 得到，因此应在该函数族之外明确列出解 $x(t)=x_2(t)$。

\textbf{已知三个不同解。}若另知解 $x_3$，令
\[
D(t)=\frac{x_3(t)-x_1(t)}{x_3(t)-x_2(t)}.
\]
则满足 $x\ne x_2$ 的解具有常交比
\[
\frac{(x-x_1)/(x-x_2)}{(x_3-x_1)/(x_3-x_2)}=C
\]
并可直接表示为
\[
\boxed{x(t)=\frac{x_1(t)-CD(t)x_2(t)}{1-CD(t)},\qquad 1-CD(t)\ne0.}
\]
取 $C=0$ 得到 $x_1$，取 $C=1$ 得到 $x_3$；解 $x_2$ 仍须在该函数族之外单独列出。三个不同解在公共连通区间上互不相交，因此定义 $D$ 所用的分母及分子均非零；各解的实际存在区间还须满足上式的分母条件。

\point{22. 特殊 Riccati 方程的初等可积指数}{special-riccati}
教材定理 2.2：$x'+a x^2=b t^m$、$a\ne0$，当
\[
m=0,-2,-\frac{4k}{2k-1},-\frac{4k}{2k+1},\qquad k=1,2,\ldots
\]
时可用反复换元化为可分离方程。$m=-2$ 时令 $z=tx$，得 $tz'=b+z-az^2$；$m=0$ 时直接分离。$ab\ne0$ 时，上列指数为教材给出的初等积分法可解情形；$b=0$ 时原方程退化为可分离方程。

\subsection{全微分方程}
\point{23. 全微分判据与区域条件}{exact}
对 $M(t,x)\dd t+N(t,x)\dd x=0$，若存在势函数 $U$ 使 $U_t=M,U_x=N$，则通积分为 $U=C$。设 $M,N\in C^1$，在矩形或单连通开区域上，充要条件为
\[
M_x=N_t.
\]
局部该条件也足够。非单连通区域上只保证闭而不保证全局恰当，必须检查闭曲线积分或分区取势函数。例 $(-x\dd t+t\dd x)/(t^2+x^2)$ 在去心平面闭而无全局单值角度势函数。

\point{24. 势函数标准公式与隐函数分支}{potential}
在矩形上选 $(t_0,x_0)$，可取
\[
U(t,x)=\int_{t_0}^t M(s,x)\dd s+\int_{x_0}^x N(t_0,u)\dd u,
\]
也可取
\[
U(t,x)=\int_{t_0}^t M(s,x_0)\dd s+\int_{x_0}^x N(t,u)\dd u.
\]
两式均满足 $U_t=M$、$U_x=N$。当 $N(t_0,x_0)\ne0$ 时，由隐函数定理，经过该点的水平集 $U=C$ 可在局部唯一表示为可微函数 $x(t)$。当 $N=0$ 时，上述隐函数定理的充分条件不成立，须根据水平集的局部形状判定是否存在可微函数分支。

\point{25. 积分因子的一般方程与两个特殊结论}{factor}
\core 设 $M,N\in C^1(D)$，在单连通开区域 $D$ 上考虑
\[
M(t,x)\dd t+N(t,x)\dd x=0.
\]
处处非零的 $C^1$ 函数 $\mu$ 为积分因子的充要条件是
\[
(\mu M)_x=(\mu N)_t,
\qquad M\mu_x-N\mu_t=(N_t-M_x)\mu.
\]

\textbf{特殊结论一：只依赖自变量 $t$ 的积分因子。}
若 $N\ne0$，且
\[
\frac{M_x(t,x)-N_t(t,x)}{N(t,x)}=h(t)
\]
为连续函数，则可取积分因子
\[
\boxed{\mu(t)=\exp\!\left(\int_{t_0}^t h(s)\dd s\right).}
\]
反之，若存在非零积分因子 $\mu=\mu(t)$，则在 $N\ne0$ 处有 $(M_x-N_t)/N=\mu'/\mu$。

\textbf{特殊结论二：只依赖未知量 $x$ 的积分因子。}
若 $M\ne0$，且
\[
\frac{N_t(t,x)-M_x(t,x)}{M(t,x)}=g(x)
\]
为连续函数，则可取积分因子
\[
\boxed{\mu(x)=\exp\!\left(\int_{x_0}^x g(u)\dd u\right).}
\]
反之，若存在非零积分因子 $\mu=\mu(x)$，则在 $M\ne0$ 处有 $(N_t-M_x)/M=\mu'/\mu$。

\textbf{适用条件。}两式中的 $t_0,x_0$ 分别取在所讨论的变量区间内；积分因子可乘任意非零常数。若分母在区域内有零点，应使用不作除法的条件
\[
M_x-N_t=N h(t)\quad\text{或}\quad N_t-M_x=M g(x)
\]
在整个讨论区域内检验；$h,g$ 须在对应变量区间连续。乘积分因子后的方程仍限于原方程定义域。

\point{27. 积分因子的非唯一性、比值与分组法}{factor-family}
若 $\mu(M\dd t+N\dd x)=\dd U$，$F$ 连续且非零，$\mu F(U)\in C^1$，则
\[
\mu F(U)(M\dd t+N\dd x)=\dd\!\left(\int^{U}F(v)\dd v\right).
\]
故 $\mu F(U)$ 仍为积分因子。两个非零 $C^1$ 积分因子的比值满足
\[
H=\mu_1/\mu_2,\qquad NH_t-MH_x=0.
\]
因此 $H$ 沿积分曲线恒定；在 $H$ 非局部常数的区域，$H=C$ 为首次积分。


\textbf{分组法的充分条件。}
记 $\omega=\omega_1+\omega_2$，$\omega_i=P_i\dd x+Q_i\dd y$，且已知
\[
\dd\Phi_1=\mu_1\omega_1,\qquad \dd\Phi_2=\mu_2\omega_2.
\]
若能找到连续单变量函数 $f,g$，使同一处处非零函数满足
\[
\boxed{\mu=\mu_1 f(\Phi_1)=\mu_2 g(\Phi_2),}
\]
且乘积具有所需的 $C^1$ 正则性，则
\[
\mu\omega=f(\Phi_1)\dd\Phi_1+g(\Phi_2)\dd\Phi_2
=\dd\!\left[\int^{\Phi_1}f(u)\dd u+\int^{\Phi_2}g(v)\dd v\right].
\]
故 $\mu$ 为 $\omega$ 的积分因子。分别存在 $\mu_1,\mu_2$ 不保证上述公共因子存在；其乘积一般不满足积分因子方程。

\textbf{可分解比值情形。}若
\[
\frac{\mu_1}{\mu_2}=A(\Phi_1)B(\Phi_2),\qquad A,B\ne0,
\]
则取 $f=1/A,g=B$，可得
\[
\boxed{\mu=\frac{\mu_1}{A(\Phi_1)}=\mu_2B(\Phi_2).}
\]
上述公式只在各函数有定义、因子非零且等式成立的区域使用。

\textbf{例：}$ (x^3y-2y^2)\dd x+x^4\dd y=0$。
在 $xy\ne0$ 上分组为
\[
\omega_1=x^3y\dd x+x^4\dd y=x^3\dd(xy),\qquad
\omega_2=-2y^2\dd x.
\]
可取 $\mu_1=x^{-3},\Phi_1=xy$ 以及 $\mu_2=-1/(2y^2),\Phi_2=x$。此时
\[
\frac{\mu_1}{\mu_2}=-\frac{2y^2}{x^3}
=(xy)^2\left(-\frac2{x^5}\right),
\]
故 $A(u)=u^2,B(v)=-2/v^5$，整体积分因子为
\[
\boxed{\mu(x,y)=\frac1{x^5y^2}.}
\]
乘回原式，得到
\[
\frac{\dd(xy)}{(xy)^2}-\frac2{x^5}\dd x
=\dd\!\left(-\frac1{xy}+\frac1{2x^4}\right)=0.
\]
通积分可写为 $1/(xy)-1/(2x^4)=C$，即
\[
y=\frac{2x^3}{1+2Cx^4},\qquad x\ne0,\quad 1+2Cx^4\ne0.
\]
构造积分因子时使用了 $xy\ne0$，故上述推导不包含零解。将 $y(x)\equiv0$ 代入原方程后，两项均为零，因此完整结果还应列出 $y(x)\equiv0$。在 $x=0$ 处，所用积分因子无定义且原方程退化；是否存在跨越该点的可微延拓，应以原方程的逐点成立为准，不能由该积分因子公式直接判定。

\subsection{隐式微分方程}
本节对应教材印刷第 42--46 页（PDF 页 52--56）。考虑一阶隐式微分方程
\[
F(t,x,x')=0.
\]
令 $p=x'(t)$，原方程成为 $F(t,x,p)=0$，并始终要求 $\dd x=p\dd t$。以下按教材的因式分解、微分法和参数法叙述；缺项方程与 Lagrange 方程的专门公式标为补充。

\point{28. 因式分解法}{implicit}
若能从 $F(t,x,p)=0$ 求出 $p$ 的全部实数解，则可将原方程化为一个或多个显式微分方程，分别求解。特别地，当 $F$ 可关于 $p$ 因式分解时，可令各因式为零，将求解转化为若干较简单的方程。

\textbf{教材例 2.4.1。}求解
\[
(x')^3-x(x')^2-t^2x'+t^2x=0.
\]
令 $p=x'$，因式分解得
\[
p^3-xp^2-t^2p+t^2x=(p-t)(p+t)(p-x)=0.
\]
于是分别求解
\[
x'=t,\qquad x'=-t,\qquad x'=x,
\]
得到三个函数族
\[
\boxed{x=\frac12t^2+C,\qquad x=-\frac12t^2+C,\qquad x=Ce^t.}
\]
这三个函数族分别满足原方程，教材将它们分别列为原方程的通解。由此可见，一阶隐式方程可以具有不同的通解表示，而不一定只有一个函数族。

\textbf{解出导数的局部条件（补充）。}若 $F\in C^1$ 且在 $F=0$ 的一点有 $F_p\ne0$，则隐函数定理保证在该点附近可将 $p$ 唯一表示为 $p=f(t,x)$。$F_p=0$ 处不满足这一充分条件，应直接分析原方程。

若两个已求出的解在某一点具有相同的函数值和一阶导数，则还可检查是否能将它们在该点连接为原方程的 $C^1$ 解。因此，若题目要求全部解，除各方程内的函数族外，还须验证可能的连接情形；连接后的函数必须逐点满足原方程。

\point{29. 微分法与缺项方程的参数解}{parameter}
\textbf{微分法。}若原方程可写为
\[
x=g(t,x'),
\]
令 $p=x'$，得到 $x=g(t,p)$。在 $p(t)$ 可微的区间内，两边对 $t$ 求导：
\[
p=g_t(t,p)+g_p(t,p)p',\qquad
\boxed{(p-g_t)\dd t=g_p\dd p.}
\]
解出 $t$ 与 $p$ 的关系后，代入 $x=g(t,p)$ 得到解。对原式求导可能增加不满足原方程的函数，所以最后须同时检验 $x=g(t,p)$ 与 $x'=p$。教材例 2.4.2 是这一方法的具体应用，见第 30 条。

\textbf{补充类型一：方程不显含未知函数 $x$。}对 $F(t,p)=0$，若可解得 $t=\phi(p)$，则
\[
\dd t=\phi'(p)\dd p,\qquad \dd x=p\phi'(p)\dd p,
\]
因此参数解为
\[
\boxed{t=\phi(p),\qquad x=\int p\phi'(p)\dd p+C.}
\]
例如教材习题 2.4 第 1(6) 题为 $(x')^3-(x')^2-2t=0$。令 $p=x'$，得到
\[
\boxed{t=\frac{p^3-p^2}{2},\qquad x=\frac38p^4-\frac13p^3+C.}
\]

\textbf{补充类型二：方程不显含自变量 $t$。}对 $F(x,p)=0$，若可解得 $x=\psi(p)$，则由 $\dd x=\psi'(p)\dd p=p\dd t$，在 $p\ne0$ 时得到
\[
\boxed{x=\psi(p),\qquad t=\int\frac{\psi'(p)}p\dd p+C.}
\]
例如教材习题 2.4 第 1(2) 题为 $(x')^3-(x')^2-x=0$，相应参数解为
\[
\boxed{x=p^3-p^2,\qquad t=\frac32p^2-2p+C.}
\]
该推导除以了 $p$，未包括 $p=x'=0$ 的情形。将 $p=0$ 代入原方程得 $x=0$，故还应列出常数解 $x(t)\equiv0$；它不能由上述参数式使 $p$ 在一个时间区间内恒为零得到。

参数式 $t=T(p)$、$x=X(p)$ 在 $T'(p)\ne0$ 处可局部解出 $p=p(t)$，从而得到可微函数 $x(t)$，并须有 $X'(p)/T'(p)=p$。在 $T'(p)=0$ 处，应根据具体参数式判断是否仍能延伸为满足原方程的可微函数，不能仅凭参数式有定义就认定该点有效。

\point{30. Clairaut 方程与包络}{clairaut}
\textbf{标准形式。}Clairaut 方程为
\[
\boxed{x=tx'+f(x').}
\]
令 $p=x'$，对 $x=tp+f(p)$ 两边求导，得
\[
p=p+tp'+f'(p)p',\qquad (t+f'(p))p'=0.
\]
由 $p'=0$ 得 $p=C$，代入原方程得到直线通解
\[
\boxed{x=Ct+f(C).}
\]
另考察 $t+f'(p)=0$。与原方程联立，得到参数曲线
\[
\boxed{t=-f'(p),\qquad x=f(p)-pf'(p).}
\]
若 $f\in C^2$ 且 $f''(p)\ne0$，则 $\dd t/\dd p=-f''(p)\ne0$，并且
\[
\frac{\dd x}{\dd p}=-pf''(p),\qquad
\frac{\dd x}{\dd t}=p.
\]
故该曲线局部给出原方程的解。在参数取值 $p=C$ 的点，它与直线 $x=Ct+f(C)$ 具有相同的函数值及斜率，构成直线族的包络。$f''=0$ 处须按具体方程检查参数曲线及相切条件。

\textbf{教材例 2.4.2。}对 $x=tx'-(x')^2/4$，有
\[
x=tp-\frac14p^2,\qquad \left(t-\frac p2\right)p'=0.
\]
由 $p=C$ 得 $x=Ct-C^2/4$；由 $p=2t$ 并代入原方程得 $x=t^2$。直线与抛物线在 $t=C/2$ 处相切，切点纵坐标为 $C^2/4$，共同斜率为 $C$。抛物线不是任意一条固定参数直线，却是原方程的解。教材在此将它称为“奇解”，并注明“见第四章”；本节使用该例与包络关系，不将第四章的一般定义另列为本次新增考点。

\point{31. 参数法与 Lagrange 方程补充}{surface-param}
\textbf{教材的一般参数法。}当原方程不易解出 $x'$，也不易写成 $x=g(t,x')$ 时，令 $p=x'$。方程 $F(t,x,p)=0$ 描述 $(t,x,p)$ 空间中的一个曲面；若其参数方程为
\[
t=T(u,v),\qquad x=X(u,v),\qquad p=P(u,v),
\]
则
\[
\dd t=T_u\dd u+T_v\dd v,\qquad
\dd x=X_u\dd u+X_v\dd v.
\]
将上述两式代入 $\dd x=p\dd t$，得到关于参数 $u,v$ 的一阶微分方程
\[
\boxed{(X_u-PT_u)\dd u+(X_v-PT_v)\dd v=0.}
\]
若解得 $u=h(v,C)$，便得到原方程的参数解
\[
t=T(h(v,C),v),\qquad x=X(h(v,C),v).
\]
若解得 $v=k(u,C)$，则相应代入 $T(u,k(u,C))$、$X(u,k(u,C))$。所得参数曲线须位于原曲面上，并在能够表示为 $x(t)$ 的部分满足 $\dd x/\dd t=P$。

\textbf{教材例 2.4.3。}对 $(x')^3-xx'+t=0$，取
\[
t=uv-v^3,\qquad x=u,\qquad p=v.
\]
由 $\dd x=p\dd t$ 得
\[
(1-v^2)\dd u+(3v^3-uv)\dd v=0.
\]
在 $|v|<1$ 上，可将它写成关于 $u(v)$ 的线性方程
\[
\frac{\dd u}{\dd v}-\frac{v}{1-v^2}u=-\frac{3v^3}{1-v^2}.
\]
乘积分因子 $\sqrt{1-v^2}$，积分得到 $u=C/\sqrt{1-v^2}+v^2+2$，因此教材在该参数区域内给出
\[
\boxed{t=\frac{Cv}{\sqrt{1-v^2}}+2v,\qquad
x=\frac{C}{\sqrt{1-v^2}}+v^2+2,\qquad |v|<1.}
\]
采用 $v$ 恢复局部函数 $x(t)$ 时，还应要求 $\dd t/\dd v\ne0$。

\textbf{本例的其他参数区域（补充）。}在 $|v|>1$ 上，同样求解关于 $u(v)$ 的线性方程，积分因子改为 $\sqrt{v^2-1}$，得到
\[
t=\frac{Cv}{\sqrt{v^2-1}}+2v,\qquad
x=\frac{C}{\sqrt{v^2-1}}+v^2+2.
\]
在 $v=1$ 或 $v=-1$ 为常值的情形，直接将 $p=\pm1$ 代入原方程，分别得到 $x=t+1$ 与 $x=1-t$；它们满足原方程，不能由上述以变化的 $v$ 为自变量的积分过程直接代替。

\textbf{补充类型：Lagrange 方程。}其具体形式为
\[
\boxed{x=t\phi(x')+\psi(x'),}
\]
其中 $\phi,\psi$ 为给定可微函数。令 $p=x'$，微分后得到
\[
p=\phi(p)+\bigl[t\phi'(p)+\psi'(p)\bigr]p'.
\]
在可取 $p$ 为自变量且 $p-\phi(p)\ne0$ 的部分，方程化为关于 $t(p)$ 的线性方程
\[
\boxed{\frac{\dd t}{\dd p}-\frac{\phi'(p)}{p-\phi(p)}t
=\frac{\psi'(p)}{p-\phi(p)}.}
\]
解出 $t=t(p,C)$ 后，代入 $x=t\phi(p)+\psi(p)$ 得到参数解。若 $p_0=\phi(p_0)$，则常数导数 $x'=p_0$ 给出原方程的直线解
\[
x(t)=p_0t+\psi(p_0).
\]
这些解因除以 $p-\phi(p)$ 被排除，故须与线性方程所得参数解一起列出。若 $\phi(p)\equiv p$，原方程即为第 30 条的 Clairaut 方程，应按该条处理，不能使用上述含零分母的线性公式。

\textbf{补充例题。}对 $x=2tx'+(x')^2$，取 $p=x'$，在 $p\ne0$ 时得到
\[
\frac{\dd t}{\dd p}+\frac2p t=-2,\qquad
\boxed{t=\frac C{p^2}-\frac23p,\qquad x=\frac{2C}{p}-\frac13p^2.}
\]
此外，$p=0$ 满足 $p=2p$，代入原方程得到解 $x(t)\equiv0$。本例为展示 Lagrange 线性化步骤的补充例题，不是教材例题编号。

\subsection{一阶微分方程组}
\point{32. 首次积分的定义与检验}{first-integral}
设 $f\in C^1(D,\R^n)$、$F\in C^1(D)$，且 $F$ 在 $D$ 的任意非空开子集上均非常数。若对系统 $x'=f(t,x)$ 的每个解 $x(t)$，$F(t,x(t))$ 为常数，则称 $F=C$ 为首次积分。其充要条件为
\[
\boxed{F_t+\sum_{j=1}^n F_{x_j}f_j=0.}
\]
上式须在整个 $D$ 上成立。$F$ 可显含 $t$，包括自治系统的首次积分。

\point{33. 独立首次积分的个数与通解}{independent-integrals}
\core 定理 2.7、2.8 及证明，见教材印刷页 54--57（PDF 页 64--67）。以下记状态向量为 $\xi\in\R^n$。

\textbf{独立性的判定。}记 $F=(F_1,\ldots,F_k)^{\mathsf T}$，$J=F_\xi$。首次积分满足 $F_t=-Jf$，因此
\[
\rank(F_t,F_\xi)=\rank\bigl(J(-f),J\bigr)=\rank J.
\]
故首次积分关于 $(t,\xi)$ 与关于 $\xi$ 的独立性等价；$n$ 维系统至多有 $n$ 个允许显含 $t$ 的函数独立首次积分。

\textbf{定理 2.7（$n$ 个独立首次积分局部确定通解）。}
设 $D\subset\R\times\R^n$ 为开集，$f\in C^1(D,\R^n)$，$F_1,\ldots,F_n\in C^1(D)$ 是 $\xi'=f(t,\xi)$ 的首次积分。若
\[
c^0=F(t_0,\xi^0),\qquad
\det J(t_0,\xi^0)=\det\frac{\partial(F_1,\ldots,F_n)}{\partial(\xi_1,\ldots,\xi_n)}(t_0,\xi^0)\ne0,
\]
则方程组 $F(t,\xi)=c$ 在 $(t_0,\xi^0,c^0)$ 附近唯一确定 $C^1$ 函数 $\xi=\phi(t,c)$。它是原系统的局部通解，且
\[
\phi_t=f(t,\phi),\qquad J(t,\phi)\phi_c=I_n,
\qquad \boxed{\det\phi_c=\frac1{\det J(t,\phi)}\ne0.}
\]

\textbf{证明。}由隐函数定理，$F(t,\phi(t,c))=c$ 在上述邻域内有唯一的 $C^1$ 解。固定 $c$，对 $t$ 求导，得到
\[
F_t(t,\phi)+J(t,\phi)\phi_t=0.
\]
另一方面，首次积分判据给出
\[
F_t(t,\phi)+J(t,\phi)f(t,\phi)=0.
\]
两式相减，得
\[
J(t,\phi)\bigl(\phi_t-f(t,\phi)\bigr)=0.
\]
缩小邻域使 $\det J\ne0$，即得 $\phi_t=f(t,\phi)$。再对 $c$ 求导，得
\[
J(t,\phi)\phi_c=I_n,\qquad \det\phi_c=(\det J(t,\phi))^{-1}\ne0,
\]
故参数独立。任一留在该邻域内的解 $\psi$ 均满足 $F(t,\psi(t))=c$；由隐函数的唯一性，$\psi(t)=\phi(t,c)$，故所得函数族构成局部通解。\hfill$\square$

该定理在 $\det J(t_0,\xi^0)\ne0$ 的点附近成立，所得 $\phi(t,c)$ 是局部单值函数。若要求在更大区域内使用同一个函数表示，还须验证隐式关系在该区域内的分支选择与唯一性。

\textbf{定理 2.8（自治系统不含时首次积分的上限）。}
设 $\xi'=f(\xi)$，$f\in C^1(U,\R^n)$。在非平衡点 $\xi^0$（即 $f(\xi^0)\ne0$）处，不含 $t$ 的 $C^1$ 首次积分中，函数独立者至多为 $n-1$ 个。更具体地，若 $F_1,\ldots,F_k$ 均为不含时首次积分，则
\[
\boxed{\rank\frac{\partial(F_1,\ldots,F_k)}{\partial(\xi_1,\ldots,\xi_n)}(\xi^0)\le n-1.}
\]

\textbf{证明（秩与零空间）。}不含时条件使判据简化为
\[
\sum_{j=1}^n\frac{\partial F_i}{\partial\xi_j}(\xi)f_j(\xi)=0
\quad(i=1,\ldots,k),\qquad J(\xi)f(\xi)=0.
\]
因 $0\ne f(\xi^0)\in\ker J(\xi^0)$，由秩与零度定理，
\[
\dim\ker J(\xi^0)\ge1,\qquad
\rank J(\xi^0)=n-\dim\ker J(\xi^0)\le n-1.
\]
\hfill$\square$

\textbf{退化情形与含时积分。}在 $f(\xi^0)=0$ 处，上述秩估计不能直接应用。反例与含时情形分别为
\[
\begin{aligned}
\xi'=0 &: &F_j(\xi)=\xi_j\ (j=1,\ldots,n),&\qquad \rank F_\xi=n;\\
x'=1,\ y'=0 &: &F_1=x-t,\ F_2=y,&\qquad \det F_{(x,y)}=1.
\end{aligned}
\]
故 $n-1$ 上限仅适用于非平衡点处的不含时首次积分。

在非平衡点附近，$n-1$ 个独立的不含时首次积分的公共水平集为一维流形，局部确定轨道的弧段。若选取满足 $f_j(\xi)\ne0$ 的坐标 $\xi_j$ 为参数，则沿该轨道恢复时间的关系为
\[
t-t_0=\int_{\xi_j(t_0)}^{\xi_j(t)}\frac{\dd u}{f_j(\xi(u))}.
\]
其中 $\xi(u)$ 表示由公共水平集确定的局部轨道参数式。由此可见，仅给定不含时积分常数一般尚未确定带时间参数的解。配套训练见第 43、61--67 题。

\point{34. 寻找首次积分的方法与降维}{reduction}
\textbf{【教材习题归纳】}以下方法由教材习题 2.5 第 2(1)--(4) 题（印刷页 57--58）归纳；标为“自拟”的例子用于补充演示，不作为教材原有的分类。设系统为 $x_i'=f_i(t,x)$，记沿解求导算子
\[
\mathcal D=\frac{\partial}{\partial t}+\sum_i f_i\frac{\partial}{\partial x_i}.
\]
求首次积分就是寻找非常数函数 $F$，使 $\mathcal DF=0$。实际计算可先寻找组合量 $U$：若 $\mathcal DU=0$，则 $U$ 已是首次积分；若 $\mathcal DU$ 只依赖于 $t,U$，则先解这个一阶方程，再将其积分常数写成 $F(t,x)$。\textbf{不必要求第一步找到的组合量立即满足导数为零。}

\textbf{（一）先求和差：寻找相同项、相反项与公共因子。}
若 $\mathcal DU=a(t)U$，则
\[
F=U\exp\left(-\int_{t_0}^{t}a(s)\,ds\right)
\]
为首次积分。例如\textbf{【自拟】} $x'=x+y,\ y'=x+y$，先取 $U=x-y$、$V=x+y$，即得
\[
U'=0,\quad V'=2V;\qquad F_1=x-y,\quad F_2=e^{-2t}(x+y).
\]
一般线性组合可设为 $U=\sum_i c_ix_i$，代入 $U'=\sum_i c_if_i$ 后选择常数 $c_i$ 使同类项相消；不应仅凭外观猜测。

\textbf{（二）有交叉乘积时，尝试平方差与带系数的平方和。}
\textbf{【教材·习题 2.5，2(2)】}对 $x'=yz,\ y'=xz,\ z'=xy$，有
\[
(x^2)'=(y^2)'=(z^2)'=2xyz.
\]
因此可取
\[
\boxed{F_1=x^2-y^2,\qquad F_2=x^2-z^2.}
\]
若各式系数不同，可设 $F=ax^2+by^2+cz^2$，由 $\mathcal DF\equiv0$ 比较系数。
例如\textbf{【自拟】} $x'=yz,\ y'=2xz,\ z'=-3xy$ 时，
\[
\mathcal DF=2(a+2b-3c)xyz,
\]
故 $a+2b-3c=0$；可选 $F_1=2x^2-y^2$、$F_2=3x^2+z^2$。\textbf{这将“凑系数”化为代数方程。}

\textbf{（三）两组合量有相同增长率时，尝试比值。}
若 $U'=aU$、$V'=aV$，其中 $a$ 可以依赖于 $t,x$，则在 $V\ne0$ 的区域
\[
\left(\frac UV\right)'=\frac{U'V-UV'}{V^2}=0.
\]
\textbf{【教材·习题 2.5，2(1)】}对
\[
x'=x^2-y^2-z^2,\qquad y'=2xy,\qquad z'=2xz,
\]
首先由 $y',z'$ 中相同的因子 $2x$ 得到 $F_1=y/z$。再取 $S=x^2+y^2+z^2$，计算得
\[
S'=2xS,\qquad z'=2xz,
\]
故在 $z\ne0$ 上得到 $F_2=S/z$。注意平方和在这里\emph{本身并不守恒}，但它与 $z$ 的比值守恒。

\textbf{（四）先有一个简单组合，再用它构造第二个积分。}
继续上题，在 $S>0$ 上有
\[
\left(\frac{x}{S}\right)'=
\frac{(x^2-y^2-z^2)S-2x^2S}{S^2}=-1,
\]
因而可取显含时间的积分
\[
\boxed{F_1=\frac yz,\qquad F_2=\frac Sz,\qquad F_3=t+\frac xS\quad(z\ne0).}
\]
上述三个积分对 $(x,y,z)$ 的 Jacobi 行列式为 $1/z^3$，故在该区域函数独立。这个过程表明，找到一个共同因子后，应继续试与它有关的比值，而不是重新任意猜一个函数。区域 $z=0$ 上的解须回到原系统求解；不能把含 $1/z$ 的公式直接代入该区域。

\textbf{（五）和差的导数不为零时，继续检查乘积。}
\textbf{【教材·习题 2.5，2(3)】}对
\[
x'=\frac{y}{x-y},\qquad y'=\frac{x}{x-y}\qquad(x\ne y),
\]
令 $U=x-y$、$V=x+y$，则
\[
U'=-1,\qquad V'=\frac VU,
\qquad (UV)'=U'V+UV'=-V+V=0.
\]
因此
\[
\boxed{F_1=t+x-y,\qquad F_2=(x-y)(x+y)=x^2-y^2.}
\]
更一般地，若非零的 $U,V$ 满足 $U'/U=a$、$V'/V=b$，可检查是否存在常数 $r,s$ 使 $ra+sb=0$；此时 $r\ln|U|+s\ln|V|$ 沿解恒定。在正值区域也可写成 $U^rV^s$。

\textbf{（六）根式或复合表达式反复出现时，将其整体换元。}
\textbf{【教材·习题 2.5，2(4)】}对
\[
x'=\frac12,\qquad y'=1+\sqrt{t-x-y},
\]
先由第一式得 $F_1=x-t/2$。再令 $u=t-x-y$，则
\[
u'=-\frac12-\sqrt u.
\]
在 $u>0$ 上分离变量，利用
\[
\int\frac{du}{\frac12+\sqrt u}
=2\sqrt u-\ln(1+2\sqrt u),
\]
得到
\[
\boxed{F_2=t+2\sqrt{t-x-y}-\ln\bigl(1+2\sqrt{t-x-y}\bigr).}
\]
这比直接在 $t,x,y$ 中猜测首次积分更有依据；边界 $u=0$ 应另行讨论。

\textbf{（七）平面旋转结构：平方和确定半径，叉乘组合确定极角。}
\textbf{【自拟】}对 $x'=-y,\ y'=x$，有
\[
(x^2+y^2)'=0,\qquad
\theta'=\frac{xy'-yx'}{x^2+y^2}=1.
\]
故在不含原点、且可选定连续极角 $\theta$ 的局部区域，可取 $F_1=x^2+y^2$、$F_2=\theta-t$。极角不能不加说明地作为整个穿孔平面上的单值函数；也可改取全局表达式
\[
\boxed{G_1=x\cos t+y\sin t,\qquad
G_2=-x\sin t+y\cos t.}
\]
二者的沿解导数均为零，且对 $(x,y)$ 的 Jacobi 行列式等于 $1$。对于 $x''+g(x)=0$ 的力学结构，应优先使用下一条的能量积分。

\textbf{检验与使用。}每得到一个候选 $F$，均应将原系统代入 $F_t+\sum_iF_{x_i}f_i$，验证其在所讨论区域恒等于零。随后检查定义域及多个积分的函数独立性；除法、对数和极角涉及的零集或分支须明确说明。对称式 $dx_i/f_i=dx_j/f_j$ 也只在相应分母非零处使用。

若 $F_{x_n}\ne0$，由隐函数定理可局部将 $F=C$ 解为 $x_n=\phi(t,x_1,\ldots,x_{n-1};C)$，代入前 $n-1$ 式得到降维系统。上述结构均不明显时，可按系统次数试低次多项式 $F$，由 $\mathcal DF\equiv0$ 比较系数；这一尝试不保证成功，也不能由未找到某种形式的积分推断系统没有首次积分。

\point{35. 自治系统的能量积分}{energy}
对 $x''+g(x)=0$，令 $v=x'$，则
\[
E=\frac12v^2+V(x),\qquad V'(x)=g(x)
\]
为首次积分。对 $x''+d(x,v)v+g(x)=0$，有
\[
E'=-d(x,v)v^2;\qquad d\ge0\Rightarrow E'\le0,\qquad d\equiv0\Rightarrow E'=0.
\]
轨道包含于能量水平集；能量常数一般不能唯一确定带时间参数的解。

\subsection{应用问题举例}
\point{36. SIR 模型的守恒量、阈值与峰值}{sir}
对教材约定 $S'=-\beta SI,I'=\beta SI-\gamma I,R'=\gamma I$，$\beta,\gamma>0$，总人数 $S+I+R=N$ 恒定。置 $\rho=\gamma/\beta$，由前两式得
\[
I+S-\rho\ln S=I_0+S_0-\rho\ln S_0,\qquad S>0.
\]
当 $I_0>0$ 时，$S$ 递减、$R$ 递增，$I'=(\beta S-\gamma)I$。若 $S_0>\rho$，未来经过 $S=\rho$ 时 $I$ 达峰值
\[
I_{\max}=I_0+S_0-\rho+\rho\ln(\rho/S_0).
\]
若 $S_0\le\rho$，初始时刻即为未来的最大值 $I_0$。教材该参数化的基本再生数为 $\mathcal R_0=\beta S_0/\gamma$；若改用发生率 $\beta SI/N$，则公式相应变成 $\beta S_0/(\gamma N)$，二者不能混用。终值满足 $I_\infty=0$ 以及 $S_\infty-\rho\ln S_\infty=I_0+S_0-\rho\ln S_0$。

\point{37. 等角轨线、正交轨线与极坐标}{trajectory}
先由曲线族 $\Phi(t,x,C)=0$ 消去 $C$ 求斜率 $H(t,x)$。与其成定向角 $\alpha$ 的轨线斜率满足
\[
x'=\frac{H+\tan\alpha}{1-H\tan\alpha},
\]
若题目给定无向交角，则须分别取定向角 $\alpha$ 与 $-\alpha$。正交轨线在 $H\ne0$ 处满足 $x'=-1/H$；斜率无定义的点应改用切向量或对称微分形式描述。

过原点的射线族的正交轨线为同心圆；与射线成固定无向交角 $0<\alpha<\pi/2$ 的轨线为对数螺线
\[
r=Ce^{\pm\theta/\tan\alpha},\qquad C>0.
\]
对极坐标曲线 $r=f(\theta)$，由 $t=r\cos\theta$、$x=r\sin\theta$ 得
\[
\frac{\dd x}{\dd t}=\frac{r'\sin\theta+r\cos\theta}{r'\cos\theta-r\sin\theta}.
\]
该式要求分母非零；当分母为零而分子非零时，曲线具有竖直切线。

\point{38. 二体问题、角动量、轨道与 Kepler 结果}{kepler}
以相对位置 $r$ 表示二体运动，标准化为 $r''=-\mu r/|r|^3$，$\mu=G(m_1+m_2)$。比角动量 $h=r\times r'$ 与比能量 $\varepsilon=|r'|^2/2-\mu/|r|$ 都守恒。$h\ne0$ 时轨道在垂直 $h$ 的平面内，且
\[
r(\theta)=\frac{p}{1+e\cos(\theta-\theta_0)},\quad
p=\frac{|h|^2}{\mu},\quad e^2=1+\frac{2\varepsilon|h|^2}{\mu^2}.
\]
$e=0$ 为圆，$0<e<1$ 为椭圆，$e=1$ 为抛物线，$e>1$ 为双曲线；$h=0$ 时为径向运动。面积速度与椭圆轨道参数满足
\[
\frac{\dd\mathcal A}{\dd t}=\frac{|h|}{2},\qquad
a=-\frac{\mu}{2\varepsilon},\qquad T^2=\frac{4\pi^2a^3}{\mu}.
\]
教材采用 $G=1$ 且固定大质量天体的近似时，以 $M$ 代替 $\mu$。

\section{第三章：只含 3.1 一般理论}
\subsection*{3.1.1 预备知识}
\point{39. 向量与矩阵的 $1$-范数及性质}{norm}
对 $x=(x_1,\ldots,x_n)^{\mathsf T}$、$A=(a_{ij})\in\R^{m\times n}$，定义
\[
\boxed{|x|_1=\sum_{i=1}^n|x_i|,\qquad
\|A\|_1=\max_{1\le j\le n}\sum_{i=1}^m|a_{ij}|.}
\]
矩阵的 $1$-范数为\textbf{各列绝对值之和的最大值}，也是向量 $1$-范数的诱导范数：
\[
\|A\|_1=\sup_{x\ne0}\frac{|Ax|_1}{|x|_1}
=\max_{|x|_1=1}|Ax|_1.
\]
\textbf{基本性质。}对同维向量 $x,y$、同型矩阵 $A,B$ 及实数 $\alpha$，有
\[
\begin{gathered}
|x|_1\ge0,\quad |x|_1=0\iff x=0,\qquad
\|A\|_1\ge0,\quad \|A\|_1=0\iff A=0,\\
|\alpha x|_1=|\alpha|\,|x|_1,\qquad
\|\alpha A\|_1=|\alpha|\,\|A\|_1,\\
|x+y|_1\le|x|_1+|y|_1,\qquad
\|A+B\|_1\le\|A\|_1+\|B\|_1.
\end{gathered}
\]
\textbf{后续估计常用：}当乘法维数相容时，
\[
\boxed{|Ax|_1\le\|A\|_1|x|_1,\qquad
\|AB\|_1\le\|A\|_1\|B\|_1,\qquad \|I_n\|_1=1.}
\]
以下省略下标时，统一将 $|x|$、$\|A\|$ 理解为上述 $1$-范数。

\point{40. 矩阵微积分与矩阵函数级数}{matrix-calculus}
对可微矩阵函数 $A,B$，其导数逐元素计算。若矩阵乘积有定义，且涉及逆矩阵时 $A$ 可逆，则有
\[
(AB)'=A'B+AB',\quad (A^{-1})'=-A^{-1}A'A^{-1},\quad
\left\|\int_a^b A(t)\dd t\right\|\le\int_a^b\|A(t)\|\dd t\quad(a\le b).
\]
矩阵积分同样逐元素计算，上述积分不等式适用于闭区间上的连续矩阵函数。有限维矩阵序列在 $1$-范数下收敛，当且仅当各矩阵元素分别收敛。对矩阵函数级数，若存在常数 $M_k\ge0$，使
\[
\|A_k(t)\|\le M_k\quad(t\in I),\qquad \sum_{k=0}^{\infty}M_k<\infty,
\]
则 $\sum_{k=0}^{\infty}A_k(t)$ 在 $I$ 上一致且绝对收敛。连续矩阵函数的一致极限连续；在闭区间上，连续项级数一致收敛时可逐项积分。逐项求导则须另加导数级数一致收敛等条件，不能仅由原级数的一致收敛推出。

\subsection*{3.1.2 解的全局存在性与唯一性}
\point{41. 连续系数线性系统的全局解}{global}
\core $A:I\to\R^{n\times n}$、$b:I\to\R^n$ 连续，$t_0\in I$，则
\[
x'=A(t)x+b(t),\quad x(t_0)=x_0
\]
在整个 $I$ 上有且仅有一个 $C^1$ 解。存在区间为系数的共同连续区间 $I$，未必为 $\R$。

\textbf{推论。}同一系统的两个解满足
\[
x_1(t_0)=x_2(t_0)\Rightarrow x_1\equiv x_2.
\]
对齐次系统，$x(t_0)=0\Rightarrow x\equiv0$；非零向量解处处非零，其单个分量可有零点。非线性方程一般无此全局结论，例如 $x'=x^2$、$x(0)=1$ 的最大存在区间为 $(-\infty,1)$。

\point{42. Picard 迭代、收敛位置与阶乘估计}{picard}
线性初值问题等价于
\[
x=x_0+\int_{t_0}^t(A(s)x(s)+b(s))\dd s.
\]
取 $x^{[0]}(t)=x_0$，定义
\[
x^{[k+1]}(t)=x_0+\int_{t_0}^t\bigl[A(s)x^{[k]}(s)+b(s)\bigr]\dd s.
\]
在任一含 $t_0$ 的紧子区间 $J\subset I$ 上，若 $\|A\|,|b|\le K$，则
\[
|x^{[k]}-x^{[k-1]}|\le(1+|x_0|)\frac{K^k|t-t_0|^k}{k!},\qquad k\ge1.
\]
故 $x^{[k]}$ 在 $J$ 上一致收敛，在 $I$ 上局部一致收敛。无界区间上的一致收敛须另行证明。两同初值解之差 $z$ 满足
\[
|z(t)|\le L\frac{(K|t-t_0|)^m}{m!}\quad(m\ge1),\qquad L=\max_{t\in J}|z(t)|;
\]
令 $m\to\infty$ 得 $z\equiv0$。

\textbf{证明复习：}教材定理 3.1，印刷页 77--80；积分等价式、阶乘估计、紧区间上一致收敛与唯一性。

\subsection*{3.1.3 齐次线性系统解的结构}
\point{43. 解空间与常系数线性组合}{solution-space}
齐次系统 $x'=A(t)x$ 的所有解构成 $n$ 维线性空间 $S(A)$。任意固定 $t_0\in I$，取值映射 $x(\cdot)\mapsto x(t_0)$ 为 $S(A)\to\R^n$ 的线性同构。任意 $n+1$ 个齐次解线性相关。

向量函数 $\phi_1,\ldots,\phi_k$ 线性无关的定义为
\[
\sum_{j=1}^k c_j\phi_j(t)\equiv0,\quad c_j\in\R
\quad\Longrightarrow\quad c_1=\cdots=c_k=0.
\]
若 $\phi_j$ 为同一连续系数齐次系统的解，则其函数线性无关等价于任一时刻的取值向量线性无关。一般向量函数不满足此等价关系；例如 $(1,0)^{\mathsf T}$、$(t,0)^{\mathsf T}$ 作为函数线性无关，而逐点线性相关。

\point{44. 解矩阵、基本解矩阵与标准化的术语}{fundamental}
$\Phi=(\phi_1,\ldots,\phi_n)$ 的每列均为系统解时叫解矩阵，满足 $\Phi'=A\Phi$。若列解线性无关，即 $\det\Phi(t)\ne0$，则为基本解矩阵。通解为 $x=\Phi c$。

\textbf{标准化约定（教材印刷页 83）。}教材称满足 $\det\Phi(t_0)=1$ 的基本解矩阵为标准基本解矩阵；本稿称满足 $\Phi(t_0)=I_n$ 者为在 $t_0$ 归一化的基本解矩阵。两条件的关系为
\[
\Phi(t_0)=I_n\Rightarrow\det\Phi(t_0)=1,\qquad
\det\Phi(t_0)=1\not\Rightarrow\Phi(t_0)=I_n.
\]

\point{45. Liouville 公式及其等价判据}{liouville}
\core 若 $W(t)=\det\Phi(t)$ 为同一齐次系统的 $n$ 个解的行列式，则
\[
W'=\tr A(t)W,\qquad
W(t)=W(t_0)\exp\int_{t_0}^t\tr A(s)\dd s.
\]
因此 $W$ 要么处处为零，要么处处非零。对同一系统的 $n$ 个解，下列条件等价：在 $I$ 上函数无关；存在一时刻 $W\ne0$；任意时刻 $W\ne0$；组成基本解矩阵。若出现孤立零点，则这些函数不可能都是同一个连续系数齐次系统的解。实系数时非零 $W$ 的符号不变；$\tr A\equiv0$ 时 $W$ 为常数。

\textbf{建议回看证明：}教材定理 3.3、3.4（印刷页 84--85），尤其“在一个时刻有相关初值＋齐次零初值唯一性”如何推出全区间相关。

\point{46. 基本解矩阵变换、反求系统与定义区间}{change-basis}
同一系统任意两个基本解矩阵满足
\[
\Psi=\Phi C,\qquad C\in\operatorname{GL}_n(\R),\quad C'=0.
\]
反之，任意常可逆矩阵 $C$ 均使 $\Phi C$ 为基本解矩阵。

给定 $C^1$ 且处处可逆的矩阵函数 $\Phi$，唯一对应的连续系数矩阵为
\[
A=\Phi'\Phi^{-1}.
\]
上述反求公式适用于 $\Phi\in C^1(I)$ 且 $\det\Phi(t)\ne0$ 对所有 $t\in I$ 成立的区间。如果 $\det\Phi$ 在某内点为零，则 $\Phi$ 不能作为包含该点的区间上的基本解矩阵。在任意 $t_0\in I$ 处，可构造归一化基本解矩阵与齐次初值解
\[
X(t)=\Phi(t)\Phi(t_0)^{-1},\quad X(t_0)=I_n,\qquad
x(t)=\Phi(t)\Phi(t_0)^{-1}x_0.
\]

\point{47. 状态转移矩阵}{transition}
令 $U(t,s)=\Phi(t)\Phi(s)^{-1}$，与选取哪一个基本解矩阵无关，且
\begin{gather*}
U(s,s)=I,\quad U(t,r)U(r,s)=U(t,s),\quad U(t,s)^{-1}=U(s,t),\\
\partial_t U(t,s)=A(t)U(t,s),\qquad
\partial_s U(t,s)=-U(t,s)A(s).
\end{gather*}
齐次初值解为 $x(t)=U(t,t_0)x_0$。出处：教材习题 3.1 第 7 题。

\subsection*{3.1.4 非齐次线性系统解的结构}
\point{48. 广义叠加与仿射解集}{nonhomogeneous}
同一非齐次系统两解之差为相应齐次解；齐次解与一个特解相加给原系统解。因此若 $x_p$ 为特解，非齐次通解为
\[
x=x_p+\Phi c,\qquad S(A,b)=x_p+S(A).
\]
$S(A,b)$ 为 $n$ 维仿射空间；当 $b\not\equiv0$ 时不是线性空间。广义叠加关系为
\[
\begin{aligned}
x_{p,j}'=Ax_{p,j}+b_j&\Rightarrow
\left(\sum_jx_{p,j}\right)'=A\sum_jx_{p,j}+\sum_jb_j,\\
x_j'=Ax_j+b,\quad\sum_jc_j=1&\Rightarrow
\left(\sum_jc_jx_j\right)'=A\sum_jc_jx_j+b
\quad(c_j\text{ 为常数}).
\end{aligned}
\]

\point{49. 常数变易法与初值公式}{variation}
\core 已知齐次基本解矩阵 $\Phi$，令 $x=\Phi c(t)$，得
\[
c'(t)=\Phi(t)^{-1}b(t),\qquad
x(t)=\Phi(t)\left[c+\int_{t_*}^t\Phi(s)^{-1}b(s)\dd s\right].
\]
初值 $x(t_0)=x_0$ 时
\[
\boxed{x(t)=\Phi(t)\Phi(t_0)^{-1}x_0+
\Phi(t)\int_{t_0}^t\Phi(s)^{-1}b(s)\dd s.}
\]
等价地，
\[
x(t)=U(t,t_0)x_0+\int_{t_0}^t U(t,s)b(s)\dd s.
\]
积分原函数的常向量可并入 $c$；矩阵乘法保持上述次序。

\point{50. 非齐次解的线性无关个数上限}{nplusone}
作业 3.1 第 10 题：$b\not\equiv0$ 时，所有非齐次解的\textbf{线性张成空间}维数恰为 $n+1$，因而至多有 $n+1$ 个线性无关的解。更具体地，取一个特解 $x_p$ 与齐次基解 $\phi_1,\ldots,\phi_n$，则 $x_p,x_p+\phi_1,\ldots,x_p+\phi_n$ 是 $n+1$ 个线性无关的原系统解。

因此
\[
\dim_{\mathrm{aff}}S(A,b)=n,\qquad
\dim\operatorname{span}S(A,b)=n+1\quad(b\not\equiv0).
\]
条件仅需 $b\not\equiv0$。证明复习：作业习题 3.1 第 10 题，增广初值 $(x(t_0),1)$ 与唯一性。

\point{51. 三角系统与组合降阶的计算结果}{triangular}
若某个分量独立，先用一阶线性公式求它，再代入其余分量逐级积分。还可找 $z=l^{\mathsf T}x$，此时 $z'=l^{\mathsf T}Ax+l^{\mathsf T}b$；若 $l^{\mathsf T}A=a(t)l^{\mathsf T}$，则得到闭合标量式 $z'=a(t)z+l^{\mathsf T}b$。

\textbf{例：}对系统 $x_1'=-2x_1/t+1$、$x_2'=(1+2/t)x_1+x_2-1$，取 $z=x_1+x_2$，则
\[
z'=z,\quad z(1)=0\Rightarrow z\equiv0,\qquad x_2=-x_1.
\]

\section*{附：全部作业与本次范围的核对}\addcontentsline{toc}{section}{附：全部作业与本次范围的核对}
\small
\begin{longtable}{p{.22\linewidth}p{.30\linewidth}p{.40\linewidth}}
\toprule 文件/版本 & 教材位置与已布置题号 & 纳入考点及处理\\\midrule\endhead
\texttt{hw1.pdf} & 习题 1.2，p.11：1(2)、2(1)、3(3)、4；习题 1.3，p.15--16：1(1)(3)、2(2) & 阶、代入验证、通解、消常数、线性判断、一阶化。原第 3(3) 题给出的函数族不是所给方程的解，不能为迎合“通解”而改题。\\
\texttt{hw2(1).pdf} & 习题 2.1，p.25：1(1)(2)(3)(5)(6)(7)、2(1)(3)、3 & 分离、平移与比值换元、初值与非唯一拼接、渐近线。\\
\texttt{hw2-2.pdf}；\texttt{hw2-2-修订.pdf} & 习题 2.2，p.33--34：1(1)(3)(5)、2、3 & 同一次作业，以修订版为主，不计成两次。积分因子、复合换元、Bernoulli 零点、线性/Riccati 两已知解。\\
\texttt{hw2-3.pdf} & 习题 2.3，p.42：1(1)(3)、2(2)(4)、3 & 全微分、积分因子、补漏零解、指定复合形式因子的充要条件。\\
\texttt{hw 10-14.pdf} & 习题 3.1：p.92 的 3(2)(3)、4(1)(3)；p.93--94 的 6 全部、8(2)、9(2)(3) & 函数无关与基本解矩阵、反求系数、定义区间、三角系统、非齐次初值公式。\\
\texttt{hw 10-25.pdf} & 习题 3.1：p.93 的 9(1)、p.94 的 10 & 常数变易、和差组合、非齐次解的线性无关上界。\\
\texttt{10.25作业2.pdf} & 习题 3.2：p.110 的 1(1)(3)(4)、2(1)(3)，p.111 的 7 & \textbf{全组超出考试截止界限，不纳入本次复习或训练。}\\\bottomrule
\end{longtable}
\normalsize
作业直接支持的主线是：2.1--2.3 的求解与条件辨析，以及 3.1 的解结构、基本解矩阵和非齐次初值。1.1 模型、2.4--2.6、3.1 的范数/迭代/状态转移矩阵当前未见单独作业文件，仍是用户指定界限内的教材内容，故保留标准结果，不能据此推断一定不考。

\textbf{少量值得回看的证明。}3.1 全局存在唯一性；Liouville 公式与“一点相关推出处处相关”；3.1 第 10 题的 $n+1$ 上界；2.3 复合积分因子判据的链式法则；Clairaut 包络的相切验证。其余条目以掌握条件和结论为主。

\section*{附：历年题型与补充结果}\addcontentsline{toc}{section}{附：历年题型与补充结果}
\textbf{仍截止 3.1（含）。}2025 年卷来自用户提供的六张试卷照片，卷首标上海交通大学致远班，2025--2026 学年第一学期；2023--2024 年文件是署名编者史瑞普的整理稿，本稿不把它直接认证为学校官方原卷。

\textbf{2025 卷题型与分值。}第 1 大题为 4 个基本求解小题，共 32 分；第 2 大题为已知特解求通解，共 20 分；第 3 大题为首次积分、独立个数及通解的确定，16 分；第 4 大题为非齐次解的 $n+1$ 上界及解集结构，10 分；第 5 大题为基本解矩阵决定系统及 Picard 误差估计，12 分；第 6 大题为地表发射的二体建模，10 分。

2023 卷也包含基本计算、Picard 迭代、Clairaut 奇解、零点非唯一性与不变号、等角轨线及能量首次积分；其一般非线性初值问题理论题可能超出当前 3.1 截止界限，本版只借鉴“先判定再论证”的问法，不扩展到第四章定理。

\subsection*{需要能独立写出的证明}
配套训练第 21--38 题专练基本证明。优先完成：积分等价式；线性系统 Picard 收敛与唯一性；齐次解空间维数；Liouville 公式（允许解矩阵奇异）；一点无关与处处无关；基本解矩阵之间只差右乘常可逆矩阵；由基本解矩阵反求系统；非齐次仿射结构与 $n+1$ 上界；常数变易公式。随后补全微分判据、复合积分因子、Bernoulli 零点与 Riccati 已知两解的推导。

上述证明分别要求明确可逆性、线性无关的常系数定义、仿射维数与线性张成空间维数，以及一致收敛的讨论区间。定理 2.7、2.8 的完整证明见第 33 条。

\subsection*{真题所问的 Picard 截断误差：标准结果}
对 $x'=A(t)x$，$x(t_0)=x_0$，$\|A(t)\|\le M$（$t\in[t_0,T]$），令 $L=T-t_0$，从 $x^{[0]}=x_0$ 起迭代，则
\[
\sup_{t\in[t_0,T]}|x(t)-x^{[n]}(t)|
\le |x_0|\sum_{k=n+1}^{\infty}\frac{(ML)^k}{k!}
\le |x_0|e^{ML}\frac{(ML)^{n+1}}{(n+1)!}.
\]
第一个不等式由相邻迭代之差的阶乘估计求尾和得到；第二个不等式由 $(n+1+j)!\ge(n+1)!j!$ 得到。此处 $n=0$ 对应常值初始近似 $x^{[0]}=x_0$。若 $M=0$ 或 $x_0=0$，则该齐次初值问题的 Picard 迭代误差为零。

\subsection*{首次积分的独立性与退化集合}
2025 卷第 3 题的系统为
\[
x'=x(y^2-z^2),\quad y'=y(z^2-x^2),\quad z'=z(x^2-y^2).
\]
可取不含时首次积分 $F_1=x^2+y^2+z^2$、$F_2=xyz$，其独立性由下式确定：
\[
\nabla F_1\times\nabla F_2
=2\bigl(x(y^2-z^2),\ y(z^2-x^2),\ z(x^2-y^2)\bigr).
\]
因此，在非平衡点处两积分函数独立；在平衡点处其 Jacobi 矩阵秩小于 $2$。该系统的平衡点集为
\[
\{(x,y,z):\text{至少两个坐标为零}\}\ \cup\ \{(x,y,z):x^2=y^2=z^2\}.
\]
两个不含时积分局部确定非平衡轨道，恢复时间参数仍需积分。由定理 2.8，非平衡点处独立的不含时首次积分至多有两个；若允许显含 $t$，独立首次积分的个数上限为三个。

教材 2.6 三个专题均纳入训练第 45--50 题：SIR 的建模、阈值与峰值/终值；等角轨线的有向角、消参数与定义域；二体问题的相对运动、角动量、能量、轨道方程与 Kepler 定律。二体的地表发射题另需写地球固定、球对称、无阻力等假设，并区分形式轨道与实际未撞地的飞行弧段。
\end{document}

